\section*{章末练习}
\begingroup\small
\begin{enumerate}[leftmargin=*]
  \item \textbf{口述概念：}无偏、一致、渐近正态和渐近有效分别约束什么？给出一个无偏但不一致的估计量。
  \item \textbf{口述概念：}区分检验大小、功效、$p$ 值、置信区间覆盖率与经济效应大小。
  \item \textbf{口述概念：}经典、异方差稳健、HAC 和簇稳健标准误各允许哪类误差结构？它们共同不能修复什么？
  \item \textbf{笔试推导：}证明 MSE 分解，并计算 $0.8\bar X$ 在 $\theta=1$、$\operatorname{Var}(\bar X)=0.25$ 时的偏差、方差和 MSE。
  \item \textbf{笔试推导：}从最小化 $\lVert y-Xb\rVert^2$ 推导正规方程，并证明残差与 $X$ 的每一列正交。
  \item \textbf{笔试推导：}推导简单回归的遗漏变量偏差公式；代入本章固定参数解释为何稳健标准误不能把斜率 $2$ 变回 $1$。

  \item \textbf{数值编程：}重复计算固定第一项与 20 项最大估计的普通正态区间，比较覆盖率。

  \item \textbf{数值编程：}逐步计算 OLS 残差与 HC0 协方差。

  \item \textbf{数值编程：}将样本标准差固定，样本量从 25 增至 100，比较标准误。
  \item \textbf{研究判断：}研究员每天查看一次同一策略的 $p$ 值，显著即停止。设计一个包含登记、停止和确认样本的修复协议。
  \item \textbf{研究判断：}加入行业控制后因子系数稳定且 HC3 显著。列出仍不能声称因果的理由与所需证据。
  \item \textbf{研究判断：}从 500 个候选因子中用 BH 选出 20 个。设计从统计发现到成本后可交易结论的证据链。
\end{enumerate}
\endgroup

\subsection*{基础衔接：独立计算}
已知标准差为 2，样本量从 25 增至 100，95\% 正态区间半宽怎样变化？
